\par\Needspace{4\baselineskip}
\begin{center}
\SourceLeaf{046}\pdfbookmark[2]{30.}{article-030}30.
\end{center}
Since none of the numbers we draw from logarithmic and trigonometric tables admits absolute precision, and each is approximate only to a certain degree, all computations performed with their aid can yield only approximately correct results. In most cases, ordinary tables accurate to the seventh decimal place, i.e. deviating from the true value by no more than half a unit in the seventh place in either direction, provide more than sufficient precision, so that the unavoidable errors are of no consequence whatever. Nevertheless, in particular cases the effect of the errors in the tables may be so greatly magnified that we are compelled to abandon an otherwise optimal method altogether and replace it with another. Such a case can also arise in the computations we have explained thus far. It will therefore be in keeping with our plan to undertake here some investigations into the degree of precision that ordinary tables allow in these computations. Although there is no room here to exhaust a subject of the greatest importance to those engaged in practical computation, we shall pursue the investigation far enough to serve our purpose and in such a way that anyone interested can refine it further and extend it to any other operations.

\par\Needspace{4\baselineskip}
\begin{center}
\hypertarget{native-046-article031}{}\pdfbookmark[2]{31.}{article-031}31.
\end{center}
Every logarithm, sine, tangent, etc. (or, more generally, every irrational quantity taken from the tables) is subject to an error that can amount to half a unit in the last digit. We shall denote this bound on the error by $\omega$, which in ordinary tables is therefore $=0{,}000\,0000\,5$. If the logarithm, etc. could not be taken directly from the tables and had to be found by interpolation, the error may be somewhat greater still for two reasons. \emph{First}, whenever the proportional part is not an integer (with the last digit taken as the unit), it is customary to use the nearest integer above or below. It is easy to see that the error can in this way be increased to almost twice its former size. We disregard this increase in the error altogether here, since that proportional part can readily have one or two decimal places \SourceLeaf{047}appended to it, and it is evident that, if the proportional part were absolutely exact, the interpolated logarithm would be subject to no greater error than the logarithms given directly in the tables, in so far as the variations of the latter may be regarded as uniform. The \emph{second} increase in error arises because that assumption is not strictly true. We also neglect it, since the effect of second and higher differences is entirely negligible in almost every case (especially if the outstanding tables whose preparation \textsc{Taylor} oversaw are used for the trigonometric quantities), and it can readily be taken into account wherever it happens to be somewhat greater. We shall therefore set the maximum unavoidable table error $=\omega$ in all cases, provided the argument (i.e. the number whose logarithm is sought, or the angle whose sine, etc., is sought) is known with absolute precision. If the argument itself is known only approximately, however, and the greatest error to which it can be subject is assumed to correspond to a variation $\omega'$ in the logarithm, etc. (a variation that can be determined by the method of differentials), the maximum error in the logarithm computed from the tables may rise to $\omega+\omega'$.

Conversely, if the argument corresponding to a given logarithm is computed with the aid of the tables, its maximum error is equal to the change in the argument corresponding to a change $\omega$ in the logarithm when the logarithm is given exactly, or to a change $\omega+\omega'$ in the logarithm when the logarithm itself may be in error by as much as $\omega'$. It will scarcely be necessary to remark that $\omega$ and $\omega'$ must be given the same sign.

If several quantities that are exact only within certain limits are added, the maximum error of the sum will be equal to the sum of the individual maximum errors, all taken with the same sign. Accordingly, in subtracting quantities that are only approximately exact, the maximum error of the difference will likewise be equal to the sum of the individual maximum errors. In multiplying or dividing a quantity that is not absolutely exact, the maximum error is increased or decreased in the same ratio as the quantity itself.
